\sectionheading{5.43 Data Structure — তথ্যের সুশৃঙ্খল বিন্যাস}

একটি Program-এ অনেকগুলো সম্পর্কিত Data আলাদা আলাদা Variable-এ রাখলে
সেগুলো পরিচালনা করা কঠিন হয়। একাধিক Data-কে নির্দিষ্ট নিয়মে
সুশৃঙ্খলভাবে সংরক্ষণ ও ব্যবহার করতে \textbf{Data Structure} ব্যবহৃত হয়।

%================================================
% Definition
%================================================

\booksubheading{Data Structure কী?}

\begin{definitionbox}

Data-কে এমনভাবে সংগঠিত ও সংরক্ষণ করার পদ্ধতি, যাতে প্রয়োজন অনুযায়ী
সহজে Access, Update ও Process করা যায়, তাকে
\textbf{Data Structure} বলে।

\end{definitionbox}

\begin{center}

\textbf{Data}
\(\longrightarrow\)
\textbf{সুশৃঙ্খলভাবে সংরক্ষণ}
\(\longrightarrow\)
\textbf{সহজে Access ও Process}

\end{center}

%================================================
% Main Structures
%================================================

\booksubheading{Python-এর গুরুত্বপূর্ণ Data Structure}

\begin{tabularx}
{\linewidth}
{
|>{\bfseries\RaggedRight\arraybackslash}p{0.22\linewidth}
|>{\RaggedRight\arraybackslash}p{0.31\linewidth}
|Y|
}

\hline
\rowcolor{softBlue}

\textbf{Data Structure}
&
\textbf{উদাহরণ}
&
\textbf{প্রধান বৈশিষ্ট্য}
\\
\hline

List
&
\textenglish{\ttfamily [10, 20, 30]}
&
পরিবর্তনযোগ্য Collection।
\\
\hline

Tuple
&
\textenglish{\ttfamily (10, 20, 30)}
&
অপরিবর্তনযোগ্য Collection।
\\
\hline

Set
&
\textenglish{\ttfamily \{10, 20, 30\}}
&
Unique Item সংরক্ষণ করে।
\\
\hline

Dictionary
&
\textenglish{\ttfamily \{"name":"Rafi"\}}
&
Key--Value Pair আকারে Data রাখে।
\\
\hline

\end{tabularx}

\begin{rememberbox}

\begin{center}

\textbf{List}
\(\longrightarrow\)
\textenglish{\ttfamily [ ]}

\qquad

\textbf{Tuple}
\(\longrightarrow\)
\textenglish{\ttfamily ( )}

\qquad

\textbf{Set}
\(\longrightarrow\)
\textenglish{\ttfamily \{ \}}

\medskip

\textbf{Dictionary}
\(\longrightarrow\)
\textenglish{\ttfamily \{key: value\}}

\end{center}

\end{rememberbox}

%================================================
% LIST
%================================================

\booksubheading{১. List}

একাধিক Item একটি Variable-এর মধ্যে ধারাবাহিকভাবে সংরক্ষণ করতে
\textbf{List} ব্যবহার করা হয়।

\begin{codeoutputbox}{List তৈরি}

\noindent
\begin{minipage}[t]{0.53\linewidth}
\codeheading

\begin{english}
\ttfamily
numbers = [10, 20, 30, 40]\\
print(numbers)
\end{english}

\end{minipage}
\hfill
\begin{minipage}[t]{0.43\linewidth}
\outputheading

\textbf{Output:}

\begin{center}
\textenglish{\ttfamily [10, 20, 30, 40]}
\end{center}

একটি List-এর মধ্যে একাধিক Item রাখা হয়েছে।

\end{minipage}

\end{codeoutputbox}

List-এর গুরুত্বপূর্ণ বৈশিষ্ট্য—

\begin{itemize}

\item Ordered।

\item Mutable বা পরিবর্তনযোগ্য।

\item Duplicate Item থাকতে পারে।

\item Index ব্যবহার করে Item Access করা যায়।

\item বিভিন্ন Data Type-এর Item রাখা যায়।

\end{itemize}

%================================================
% Index
%================================================

\booksubheading{List-এর Index}

Python-এ Index সাধারণত
\textenglish{\ttfamily 0}
থেকে শুরু হয়।

\begin{codeoutputbox}{Index ব্যবহার করে Item Access}

\noindent
\begin{minipage}[t]{0.53\linewidth}
\codeheading

\begin{english}
\ttfamily
numbers = [10, 20, 30, 40]\\
print(numbers[0])\\
print(numbers[2])
\end{english}

\end{minipage}
\hfill
\begin{minipage}[t]{0.43\linewidth}
\outputheading

\begin{english}
\ttfamily
10\\
30
\end{english}

\medskip

প্রথম Item-এর Index
\textenglish{\ttfamily 0}
এবং তৃতীয় Item-এর Index
\textenglish{\ttfamily 2}।

\end{minipage}

\end{codeoutputbox}

\begin{tabularx}
{\linewidth}
{
|>{\bfseries\centering\arraybackslash}p{0.20\linewidth}
|>{\centering\arraybackslash}Y
|>{\centering\arraybackslash}Y
|>{\centering\arraybackslash}Y
|>{\centering\arraybackslash}Y|
}

\hline
\rowcolor{softBlue}

\textbf{Item}
&
10
&
20
&
30
&
40
\\
\hline

\textbf{Index}
&
0
&
1
&
2
&
3
\\
\hline

\end{tabularx}

\begin{rememberbox}

\textbf{\color{deepOrange}MCQ Trap:}
Python-এর প্রথম Item-এর Index সাধারণত
\textenglish{\ttfamily 0}।

\end{rememberbox}

%================================================
% Negative Index
%================================================

\booksubheading{Negative Index}

Negative Index ব্যবহার করে শেষের দিক থেকে Item Access করা যায়।

\begin{codeoutputbox}{Negative Index}

\noindent
\begin{minipage}[t]{0.53\linewidth}
\codeheading

\begin{english}
\ttfamily
numbers = [10, 20, 30, 40]\\
print(numbers[-1])\\
print(numbers[-2])
\end{english}

\end{minipage}
\hfill
\begin{minipage}[t]{0.43\linewidth}
\outputheading

\begin{english}
\ttfamily
40\\
30
\end{english}

\medskip

\textenglish{\ttfamily -1}
\(\longrightarrow\)
শেষ Item

\textenglish{\ttfamily -2}
\(\longrightarrow\)
শেষের আগের Item

\end{minipage}

\end{codeoutputbox}

%================================================
% List Update
%================================================

\booksubheading{List-এর Item পরিবর্তন}

List \textbf{Mutable}, তাই তৈরির পরে Item পরিবর্তন করা যায়।

\begin{codeoutputbox}{Item Update}

\noindent
\begin{minipage}[t]{0.53\linewidth}
\codeheading

\begin{english}
\ttfamily
numbers = [10, 20, 30]\\
numbers[1] = 50\\
print(numbers)
\end{english}

\end{minipage}
\hfill
\begin{minipage}[t]{0.43\linewidth}
\outputheading

\begin{center}
\textenglish{\ttfamily [10, 50, 30]}
\end{center}

Index
\textenglish{\ttfamily 1}
-এর Value
\textenglish{\ttfamily 20}
থেকে
\textenglish{\ttfamily 50}
হয়েছে।

\end{minipage}

\end{codeoutputbox}

%================================================
% append insert
%================================================

\booksubheading{List-এ Item যোগ}

\begin{codeoutputbox}{\textenglish{\ttfamily append()} ও \textenglish{\ttfamily insert()}}

\noindent
\begin{minipage}[t]{0.53\linewidth}
\codeheading

\begin{english}
\ttfamily
numbers = [10, 20]\\
numbers.append(30)\\
numbers.insert(1, 15)\\
print(numbers)
\end{english}

\end{minipage}
\hfill
\begin{minipage}[t]{0.43\linewidth}
\outputheading

\begin{center}
\textenglish{\ttfamily [10, 15, 20, 30]}
\end{center}

\textenglish{\ttfamily append(30)}
শেষে 30 যোগ করেছে।

\textenglish{\ttfamily insert(1, 15)}
Index 1-এ 15 যোগ করেছে।

\end{minipage}

\end{codeoutputbox}

%================================================
% remove pop
%================================================

\booksubheading{List থেকে Item বাদ}

\begin{codeoutputbox}{\textenglish{\ttfamily remove()} ও \textenglish{\ttfamily pop()}}

\noindent
\begin{minipage}[t]{0.54\linewidth}
\codeheading

\begin{english}
\ttfamily
numbers = [10, 20, 30, 40]\\
numbers.remove(20)\\
x = numbers.pop(1)\\
print(x)\\
print(numbers)
\end{english}

\end{minipage}
\hfill
\begin{minipage}[t]{0.42\linewidth}
\outputheading

\begin{english}
\ttfamily
30\\{}
[10, 40]
\end{english}

\medskip

\textenglish{\ttfamily remove(20)}
Value 20 বাদ দেয়।

\textenglish{\ttfamily pop(1)}
Index 1-এর Item সরিয়ে সেটি Return করে।

\end{minipage}

\end{codeoutputbox}

%================================================
% Slicing
%================================================

\booksubheading{List Slicing}

List-এর নির্দিষ্ট অংশ নেওয়াকে \textbf{Slicing} বলা হয়।

\begin{codeoutputbox}{Slicing}

\noindent
\begin{minipage}[t]{0.53\linewidth}
\codeheading

\begin{english}
\ttfamily
numbers = [10, 20, 30, 40, 50]\\
print(numbers[1:4])
\end{english}

\end{minipage}
\hfill
\begin{minipage}[t]{0.43\linewidth}
\outputheading

\begin{center}
\textenglish{\ttfamily [20, 30, 40]}
\end{center}

\textenglish{\ttfamily [1:4]}
এ Index 1 থেকে শুরু হলেও Index 4 অন্তর্ভুক্ত হয় না।

\medskip

\textbf{Start Included, Stop Excluded}

\end{minipage}

\end{codeoutputbox}

%================================================
% List Loop
%================================================

\booksubheading{List-এর উপর Loop}

\begin{codeoutputbox}{List Traversal}

\noindent
\begin{minipage}[t]{0.53\linewidth}
\codeheading

\begin{english}
\ttfamily
numbers = [10, 20, 30]\\
for x in numbers:\\
\hspace*{4mm}print(x)
\end{english}

\end{minipage}
\hfill
\begin{minipage}[t]{0.43\linewidth}
\outputheading

\begin{english}
\ttfamily
10\\
20\\
30
\end{english}

\medskip

প্রতিটি Iteration-এ List-এর একটি করে Item
\textenglish{\ttfamily x}
-এর মধ্যে আসে।

\end{minipage}

\end{codeoutputbox}

%================================================
% TUPLE
%================================================

\booksubheading{২. Tuple}

List-এর মতো Tuple-এও একাধিক Item রাখা যায়। তবে Tuple-এর Item
তৈরির পরে সাধারণভাবে পরিবর্তন করা যায় না।

\begin{codeoutputbox}{Tuple তৈরি ও Access}

\noindent
\begin{minipage}[t]{0.53\linewidth}
\codeheading

\begin{english}
\ttfamily
numbers = (10, 20, 30)\\
print(numbers[1])
\end{english}

\end{minipage}
\hfill
\begin{minipage}[t]{0.43\linewidth}
\outputheading

\begin{center}
\textenglish{\ttfamily 20}
\end{center}

Tuple Ordered এবং Index ব্যবহার করে Item Access করা যায়।

কিন্তু এটি \textbf{Immutable}।

\end{minipage}

\end{codeoutputbox}

Tuple-এর গুরুত্বপূর্ণ বৈশিষ্ট্য—

\begin{itemize}

\item Ordered।

\item Immutable বা অপরিবর্তনযোগ্য।

\item Duplicate Item থাকতে পারে।

\item Index ব্যবহার করা যায়।

\end{itemize}

\begin{rememberbox}

\begin{center}

\textbf{List = Mutable}
\qquad
\textbf{Tuple = Immutable}

\end{center}

\end{rememberbox}

%================================================
% Single Item Tuple
%================================================

\booksubheading{Single Item Tuple}

একটি Item-এর Tuple তৈরি করতে Comma গুরুত্বপূর্ণ।

\begin{codeoutputbox}{Comma-এর গুরুত্ব}

\noindent
\begin{minipage}[t]{0.48\linewidth}
\codeheading

\begin{english}
\ttfamily
x = (10,)\\
y = (10)
\end{english}

\end{minipage}
\hfill
\begin{minipage}[t]{0.48\linewidth}
\outputheading

\textenglish{\ttfamily x = (10,)}
\(\longrightarrow\)
Tuple

\medskip

\textenglish{\ttfamily y = (10)}
\(\longrightarrow\)
সাধারণ Integer Value

\end{minipage}

\end{codeoutputbox}

\begin{rememberbox}

\textbf{\color{deepOrange}MCQ Trap:}
Single Item Tuple তৈরি করতে Item-এর পরে Comma প্রয়োজন।

\begin{center}
\textenglish{\ttfamily (10,)}
\end{center}

\end{rememberbox}

%================================================
% SET
%================================================

\booksubheading{৩. Set}

Unique Item-এর Collection সংরক্ষণ করতে \textbf{Set} ব্যবহার করা হয়।

Set-এর গুরুত্বপূর্ণ বৈশিষ্ট্য—

\begin{compacttwocolpoints}

\item Duplicate Item রাখে না।

\item Indexing সমর্থন করে না।

\item Item-এর নির্দিষ্ট Positional Order-এর উপর নির্ভর করা উচিত নয়।

\item Mutable।

\item Membership Testing-এ উপযোগী।

\end{compacttwocolpoints}

%================================================
% Set Duplicate
%================================================

\begin{codeoutputbox}{Set-এ Duplicate Item}

\noindent
\begin{minipage}[t]{0.53\linewidth}
\codeheading

\begin{english}
\ttfamily
numbers = \{10, 20, 20, 30\}\\
print(numbers)
\end{english}

\end{minipage}
\hfill
\begin{minipage}[t]{0.43\linewidth}
\outputheading

Set-এ
\textenglish{\ttfamily 20}
একবারই থাকবে।

উদাহরণস্বরূপ—

\begin{center}
\textenglish{\ttfamily \{10, 20, 30\}}
\end{center}

Item প্রদর্শনের Order ভিন্ন হতে পারে।

\end{minipage}

\end{codeoutputbox}

\begin{rememberbox}

\begin{center}

\textbf{Set}
\(\longrightarrow\)
\textbf{Duplicate Item বাদ দেয়}

\end{center}

\end{rememberbox}

%================================================
% Set Add Remove
%================================================

\booksubheading{Set-এ Item যোগ ও বাদ}

\begin{codeoutputbox}{\textenglish{\ttfamily add()} ও \textenglish{\ttfamily remove()}}

\noindent
\begin{minipage}[t]{0.53\linewidth}
\codeheading

\begin{english}
\ttfamily
numbers = \{10, 20\}\\
numbers.add(30)\\
numbers.remove(20)\\
print(numbers)
\end{english}

\end{minipage}
\hfill
\begin{minipage}[t]{0.43\linewidth}
\outputheading

Result-এ 30 যোগ হবে এবং 20 বাদ যাবে।

উদাহরণস্বরূপ—

\begin{center}
\textenglish{\ttfamily \{10, 30\}}
\end{center}

\textenglish{\ttfamily add()}
Item যোগ করে,
\textenglish{\ttfamily remove()}
Item বাদ দেয়।

\end{minipage}

\end{codeoutputbox}

%================================================
% Empty Set
%================================================

\booksubheading{Empty Set}

\begin{codeoutputbox}{\textenglish{\ttfamily \{\}} বনাম \textenglish{\ttfamily set()}}

\noindent
\begin{minipage}[t]{0.48\linewidth}
\codeheading

\begin{english}
\ttfamily
a = \{\}\\
b = set()
\end{english}

\end{minipage}
\hfill
\begin{minipage}[t]{0.48\linewidth}
\outputheading

\textenglish{\ttfamily a = \{\}}
\(\longrightarrow\)
Empty Dictionary

\medskip

\textenglish{\ttfamily b = set()}
\(\longrightarrow\)
Empty Set

\end{minipage}

\end{codeoutputbox}

\begin{rememberbox}

\textbf{\color{deepOrange}MCQ Trap}

\begin{center}

\textenglish{\ttfamily \{\}}
\(\longrightarrow\)
\textbf{Empty Dictionary}

\qquad

\textenglish{\ttfamily set()}
\(\longrightarrow\)
\textbf{Empty Set}

\end{center}

\end{rememberbox}

%================================================
% DICTIONARY
%================================================

\booksubheading{৪. Dictionary}

Data-কে \textbf{Key--Value Pair} আকারে সংরক্ষণ করতে
Dictionary ব্যবহার করা হয়।

\begin{codeoutputbox}{Dictionary তৈরি}

\noindent
\begin{minipage}[t]{0.56\linewidth}
\codeheading

\begin{english}
\ttfamily
student = \{\\
\hspace*{4mm}"name": "Rafi",\\
\hspace*{4mm}"age": 18,\\
\hspace*{4mm}"marks": 85\\
\}
\end{english}

\end{minipage}
\hfill
\begin{minipage}[t]{0.40\linewidth}
\outputheading

এখানে—

\textenglish{\ttfamily "name"}
\(\longrightarrow\)
Key

\textenglish{\ttfamily "Rafi"}
\(\longrightarrow\)
Value

\medskip

Dictionary-তে Data
\textbf{Key--Value Pair}
আকারে থাকে।

\end{minipage}

\end{codeoutputbox}

%================================================
% Dict access
%================================================

\booksubheading{Dictionary থেকে Value Access}

\begin{codeoutputbox}{Key ব্যবহার করে Access}

\noindent
\begin{minipage}[t]{0.56\linewidth}
\codeheading

\begin{english}
\ttfamily
student = \{\\
\hspace*{4mm}"name": "Rafi",\\
\hspace*{4mm}"age": 18\\
\}\\
print(student["name"])
\end{english}

\end{minipage}
\hfill
\begin{minipage}[t]{0.40\linewidth}
\outputheading

\begin{center}
\textenglish{\ttfamily Rafi}
\end{center}

Dictionary-তে সাধারণত Position নয়,
\textbf{Key} ব্যবহার করে Value Access করা হয়।

\end{minipage}

\end{codeoutputbox}

%================================================
% Dict update add
%================================================

\booksubheading{Dictionary-এর Value পরিবর্তন ও নতুন Item যোগ}

\begin{codeoutputbox}{Update ও Add}

\noindent
\begin{minipage}[t]{0.56\linewidth}
\codeheading

\begin{english}
\ttfamily
student = \{\\
\hspace*{4mm}"name": "Rafi",\\
\hspace*{4mm}"age": 18\\
\}\\
student["age"] = 19\\
student["city"] = "Rajshahi"
\end{english}

\end{minipage}
\hfill
\begin{minipage}[t]{0.40\linewidth}
\outputheading

\textenglish{\ttfamily "age"}
-এর Value 18 থেকে 19 হয়েছে।

\medskip

\textenglish{\ttfamily "city"}
নামে নতুন Key--Value Pair যোগ হয়েছে।

\medskip

Dictionary \textbf{Mutable}।

\end{minipage}

\end{codeoutputbox}

%================================================
% Dict pop
%================================================

\booksubheading{Dictionary থেকে Item বাদ}

\begin{codeoutputbox}{\textenglish{\ttfamily pop()}}

\noindent
\begin{minipage}[t]{0.54\linewidth}
\codeheading

\begin{english}
\ttfamily
student = \{"name": "Rafi", "age": 18\}\\
student.pop("age")\\
print(student)
\end{english}

\end{minipage}
\hfill
\begin{minipage}[t]{0.42\linewidth}
\outputheading

\begin{center}
\textenglish{\ttfamily \{"name": "Rafi"\}}
\end{center}

\textenglish{\ttfamily pop("age")}
দিয়ে
\textenglish{\ttfamily "age"}
Key ও তার Value বাদ গেছে।

\end{minipage}

\end{codeoutputbox}

%================================================
% Dictionary Loop
%================================================

\booksubheading{Dictionary-এর উপর Loop}

\begin{codeoutputbox}{সরাসরি Dictionary Iteration}

\noindent
\begin{minipage}[t]{0.55\linewidth}
\codeheading

\begin{english}
\ttfamily
student = \{"name": "Rafi", "age": 18\}\\
for x in student:\\
\hspace*{4mm}print(x)
\end{english}

\end{minipage}
\hfill
\begin{minipage}[t]{0.41\linewidth}
\outputheading

\begin{english}
\ttfamily
name\\
age
\end{english}

\medskip

Dictionary-এর উপর সরাসরি Loop চালালে
\textbf{Key} পাওয়া যায়।

\end{minipage}

\end{codeoutputbox}

\begin{rememberbox}

\textbf{\color{deepOrange}MCQ Trap:}

\begin{center}

\textbf{Dictionary-এর উপর সরাসরি Loop}
\(\longrightarrow\)
\textbf{Key-এর উপর Iteration}

\end{center}

\end{rememberbox}

%================================================
% items
%================================================

\booksubheading{Key ও Value একসঙ্গে পাওয়া}

\begin{codeoutputbox}{\textenglish{\ttfamily items()} Method}

\noindent
\begin{minipage}[t]{0.56\linewidth}
\codeheading

\begin{english}
\ttfamily
student = \{"name": "Rafi", "age": 18\}\\
for key, value in student.items():\\
\hspace*{4mm}print(key, value)
\end{english}

\end{minipage}
\hfill
\begin{minipage}[t]{0.40\linewidth}
\outputheading

\begin{english}
\ttfamily
name Rafi\\
age 18
\end{english}

\medskip

\textenglish{\ttfamily items()}
Key ও Value একসঙ্গে দেয়।

\end{minipage}

\end{codeoutputbox}

%================================================
% Dictionary methods
%================================================

\booksubheading{Dictionary-এর গুরুত্বপূর্ণ Method}

\begin{tabularx}
{\linewidth}
{
|>{\bfseries\RaggedRight\arraybackslash}p{0.30\linewidth}
|Y|
}

\hline
\rowcolor{softBlue}

\textbf{Method}
&
\textbf{কাজ}
\\
\hline

\textenglish{\ttfamily keys()}
&
সব Key পাওয়া যায়।
\\
\hline

\textenglish{\ttfamily values()}
&
সব Value পাওয়া যায়।
\\
\hline

\textenglish{\ttfamily items()}
&
Key--Value Pair পাওয়া যায়।
\\
\hline

\textenglish{\ttfamily get()}
&
Key অনুযায়ী Value পাওয়া যায়।
\\
\hline

\textenglish{\ttfamily pop()}
&
নির্দিষ্ট Key-এর Item বাদ দেয়।
\\
\hline

\end{tabularx}

\begin{codeoutputbox}{\textenglish{\ttfamily get()} Example}

\noindent
\begin{minipage}[t]{0.55\linewidth}
\codeheading

\begin{english}
\ttfamily
student = \{"name": "Rafi", "age": 18\}\\
print(student.get("name"))
\end{english}

\end{minipage}
\hfill
\begin{minipage}[t]{0.41\linewidth}
\outputheading

\begin{center}
\textenglish{\ttfamily Rafi}
\end{center}

\textenglish{\ttfamily get("name")}
Key অনুযায়ী Value Return করেছে।

\end{minipage}

\end{codeoutputbox}

%================================================
% len
%================================================

\booksubheading{\textenglish{\ttfamily len()} Function}

List, Tuple, Set ও Dictionary-তে Item সংখ্যা নির্ণয়ের জন্য
\textenglish{\ttfamily len()}
ব্যবহার করা যায়।

\begin{codeoutputbox}{\textenglish{\ttfamily len()} Example}

\noindent
\begin{minipage}[t]{0.53\linewidth}
\codeheading

\begin{english}
\ttfamily
numbers = [10, 20, 30]\\
print(len(numbers))
\end{english}

\end{minipage}
\hfill
\begin{minipage}[t]{0.43\linewidth}
\outputheading

\begin{center}
\textenglish{\ttfamily 3}
\end{center}

List-এ মোট 3টি Item আছে।

Dictionary-এর ক্ষেত্রে
\textenglish{\ttfamily len()}
Key--Value Pair-এর সংখ্যা নির্দেশ করে।

\end{minipage}

\end{codeoutputbox}

%================================================
% Membership
%================================================

\booksubheading{Membership Operator}

কোনো Item Collection-এর মধ্যে আছে কি না পরীক্ষা করতে
\textenglish{\ttfamily in}
ব্যবহার করা যায়।

\begin{codeoutputbox}{List-এ Membership Test}

\noindent
\begin{minipage}[t]{0.53\linewidth}
\codeheading

\begin{english}
\ttfamily
numbers = [10, 20, 30]\\
print(20 in numbers)
\end{english}

\end{minipage}
\hfill
\begin{minipage}[t]{0.43\linewidth}
\outputheading

\begin{center}
\textenglish{\ttfamily True}
\end{center}

20 List-এর মধ্যে আছে, তাই ফলাফল
\textenglish{\ttfamily True}।

\end{minipage}

\end{codeoutputbox}

\begin{codeoutputbox}{Dictionary-তে Membership Test}

\noindent
\begin{minipage}[t]{0.56\linewidth}
\codeheading

\begin{english}
\ttfamily
student = \{"name": "Rafi", "age": 18\}\\
print("name" in student)
\end{english}

\end{minipage}
\hfill
\begin{minipage}[t]{0.40\linewidth}
\outputheading

\begin{center}
\textenglish{\ttfamily True}
\end{center}

Dictionary-এর ক্ষেত্রে
\textenglish{\ttfamily in}
সাধারণত \textbf{Key} পরীক্ষা করে।

\end{minipage}

\end{codeoutputbox}

%================================================
% Main Comparison
%================================================

\booksubheading{List, Tuple, Set ও Dictionary-এর পার্থক্য}

\begin{tabularx}
{\linewidth}
{
|>{\bfseries\RaggedRight\arraybackslash}p{0.19\linewidth}
|>{\centering\arraybackslash}Y
|>{\centering\arraybackslash}Y
|>{\centering\arraybackslash}Y
|>{\centering\arraybackslash}Y|
}

\hline
\rowcolor{softBlue}

\textbf{বৈশিষ্ট্য}
&
\textbf{List}
&
\textbf{Tuple}
&
\textbf{Set}
&
\textbf{Dictionary}
\\
\hline

চিহ্ন
&
\textenglish{\ttfamily [ ]}
&
\textenglish{\ttfamily ( )}
&
\textenglish{\ttfamily \{ \}}
&
\textenglish{\ttfamily \{key:value\}}
\\
\hline

Ordered
&
হ্যাঁ
&
হ্যাঁ
&
Positional নয়
&
হ্যাঁ
\\
\hline

Mutable
&
হ্যাঁ
&
না
&
হ্যাঁ
&
হ্যাঁ
\\
\hline

Duplicate
&
হ্যাঁ
&
হ্যাঁ
&
না
&
Key নয়
\\
\hline

Access
&
Index
&
Index
&
Index নয়
&
Key
\\
\hline

প্রধান ব্যবহার
&
পরিবর্তনযোগ্য Data
&
স্থির Data
&
Unique Item
&
Key--Value Data
\\
\hline

\end{tabularx}

%================================================
% Mutable Immutable
%================================================

\booksubheading{Mutable বনাম Immutable}

\begin{tabularx}
{\linewidth}
{
|>{\bfseries\RaggedRight\arraybackslash}p{0.27\linewidth}
|>{\RaggedRight\arraybackslash}p{0.32\linewidth}
|Y|
}

\hline
\rowcolor{softBlue}

\textbf{ধরন}
&
\textbf{অর্থ}
&
\textbf{উদাহরণ}
\\
\hline

Mutable
&
তৈরির পরে Content পরিবর্তন করা যায়।
&
List, Set, Dictionary
\\
\hline

Immutable
&
তৈরির পরে Content পরিবর্তন করা যায় না।
&
Tuple
\\
\hline

\end{tabularx}

\begin{rememberbox}

\begin{center}

\textbf{List + Set + Dictionary}
\(\longrightarrow\)
\textbf{Mutable}

\qquad

\textbf{Tuple}
\(\longrightarrow\)
\textbf{Immutable}

\end{center}

\end{rememberbox}

%================================================
% Nested Data Structure
%================================================

\booksubheading{Nested Data Structure}

একটি Data Structure-এর ভিতরে অন্য Data Structure রাখা যায়।

\begin{codeoutputbox}{List-এর ভিতরে List}

\noindent
\begin{minipage}[t]{0.56\linewidth}
\codeheading

\begin{english}
\ttfamily
students = [\\
\hspace*{4mm}["Rafi", 85],\\
\hspace*{4mm}["Nila", 90]\\
]\\
print(students[0][0])
\end{english}

\end{minipage}
\hfill
\begin{minipage}[t]{0.40\linewidth}
\outputheading

\begin{center}
\textenglish{\ttfamily Rafi}
\end{center}

\textenglish{\ttfamily students[0]}
প্রথম Inner List।

তার
\textenglish{\ttfamily [0]}
Index-এর Item হলো
\textenglish{\ttfamily "Rafi"}।

\end{minipage}

\end{codeoutputbox}

%================================================
% Dictionary with list
%================================================

\begin{codeoutputbox}{Dictionary-এর ভিতরে List}

\noindent
\begin{minipage}[t]{0.56\linewidth}
\codeheading

\begin{english}
\ttfamily
student = \{\\
\hspace*{4mm}"name": "Rafi",\\
\hspace*{4mm}"marks": [80, 85, 90]\\
\}\\
print(student["marks"][1])
\end{english}

\end{minipage}
\hfill
\begin{minipage}[t]{0.40\linewidth}
\outputheading

\begin{center}
\textenglish{\ttfamily 85}
\end{center}

প্রথমে
\textenglish{\ttfamily "marks"}
Key-এর List পাওয়া হয়েছে।

তারপর Index
\textenglish{\ttfamily 1}
থেকে 85 পাওয়া গেছে।

\end{minipage}

\end{codeoutputbox}

%================================================
% When to use
%================================================

\chFiveInput{topics/43_data_structures/python_extensions.tex}

\booksubheading{কোন Data Structure কখন ব্যবহার করব?}

\begin{tabularx}
{\linewidth}
{
|>{\bfseries\RaggedRight\arraybackslash}p{0.58\linewidth}
|Y|
}

\hline
\rowcolor{softBlue}

\textbf{প্রয়োজন}
&
\textbf{উপযোগী}
\\
\hline

Item পরিবর্তন করতে হবে
&
List
\\
\hline

Data অপরিবর্তিত রাখতে চাই
&
Tuple
\\
\hline

Duplicate Item বাদ দিতে চাই
&
Set
\\
\hline

Unique Collection প্রয়োজন
&
Set
\\
\hline

Key দিয়ে Data Access করতে চাই
&
Dictionary
\\
\hline

নাম--Value / ID--Value ধরনের Data
&
Dictionary
\\
\hline

\end{tabularx}

%================================================
% Benefits
%================================================

\booksubheading{Data Structure ব্যবহারের সুবিধা}

\begin{compacttwocolpoints}

\item অনেক Data একসঙ্গে সংরক্ষণ করা যায়।

\item সম্পর্কিত Data সুশৃঙ্খলভাবে রাখা যায়।

\item Data সহজে Access ও Process করা যায়।

\item Loop ব্যবহার করে Collection-এর Item নিয়ে কাজ করা যায়।

\item Search, Update ও Delete সহজ হয়।

\item Program আরও সংগঠিত ও পরিচালনাযোগ্য হয়।

\end{compacttwocolpoints}

%================================================
% MCQ
%================================================

\booksubheading{গুরুত্বপূর্ণ MCQ তথ্য}

\begin{rememberbox}

\begin{tabularx}
{\linewidth}
{
>{\bfseries\RaggedRight\arraybackslash}p{0.37\linewidth}
Y
}

List
&
\textenglish{\ttfamily [ ]}; Mutable; Duplicate থাকতে পারে।
\\[0.45mm]

Tuple
&
\textenglish{\ttfamily ( )}; Immutable; Duplicate থাকতে পারে।
\\[0.45mm]

Set
&
Unique Item; Indexing নেই।
\\[0.45mm]

Dictionary
&
Key--Value Pair; Key দিয়ে Access।
\\[0.45mm]

প্রথম Index
&
\textenglish{\ttfamily 0}
\\[0.45mm]

শেষ Item
&
Negative Index
\textenglish{\ttfamily -1}
\\[0.45mm]

Empty Dictionary
&
\textenglish{\ttfamily \{\}}
\\[0.45mm]

Empty Set
&
\textenglish{\ttfamily set()}
\\[0.45mm]

Single Item Tuple
&
Comma প্রয়োজন:
\textenglish{\ttfamily (10,)}
\\[0.45mm]

Dictionary Loop
&
সরাসরি Loop করলে Key পাওয়া যায়।
\\[0.45mm]

Dictionary-তে \textenglish{\ttfamily in}
&
সাধারণত Key পরীক্ষা করে।
\\[0.45mm]

\textenglish{\ttfamily len()}
&
Item / Entry-এর সংখ্যা নির্ণয় করে।
\\

\end{tabularx}

\end{rememberbox}

%================================================
% Common Mistakes
%================================================

\booksubheading{সাধারণ ভুল}

\begin{rememberbox}

\begin{itemize}

\item
\textenglish{\ttfamily \{\}}
-কে Empty Set মনে করা।

\item Set-এ Index ব্যবহার করার চেষ্টা করা।

\item Tuple-এর Item পরিবর্তনের চেষ্টা করা।

\item List ও Tuple-এর Mutable/Immutable বৈশিষ্ট্য গুলিয়ে ফেলা।

\item Dictionary-তে Index দিয়ে Value Access করার চেষ্টা করা।

\item Dictionary-এর উপর সরাসরি Loop চালালে Value পাওয়া যাবে মনে করা।

\item Single Item Tuple-এ Comma না দেওয়া।

\item List Slicing-এ Stop Index Included মনে করা।

\item Set-এর Item সবসময় একই Position/Order-এ থাকবে ধরে নেওয়া।

\end{itemize}

\end{rememberbox}

%================================================
% Final Reminder
%================================================

\begin{rememberbox}

\textbf{\color{deepOrange}মনে রাখবে}

\begin{tabularx}
{\linewidth}
{
>{\bfseries\RaggedRight\arraybackslash}p{0.23\linewidth}
Y
>{\bfseries\RaggedRight\arraybackslash}p{0.23\linewidth}
Y
}

List
&
Ordered + Mutable
&
Tuple
&
Ordered + Immutable
\\[0.5mm]

Set
&
Unique Item
&
Dictionary
&
Key--Value Pair
\\[0.5mm]

\textenglish{\ttfamily [ ]}
&
List
&
\textenglish{\ttfamily ( )}
&
Tuple
\\[0.5mm]

\textenglish{\ttfamily \{10,20\}}
&
Set
&
\textenglish{\ttfamily \{"name":"Rafi"\}}
&
Dictionary
\\[0.5mm]

\textenglish{\ttfamily \{\}}
&
Empty Dictionary
&
\textenglish{\ttfamily set()}
&
Empty Set
\\[0.5mm]

List / Tuple
&
Index দিয়ে Access
&
Dictionary
&
Key দিয়ে Access
\\[0.5mm]

\textenglish{\ttfamily len()}
&
Item Count
&
\textenglish{\ttfamily in}
&
Membership Test
\\

\end{tabularx}

\medskip

\begin{center}

\textbf{পরিবর্তন করতে হবে}
\(\longrightarrow\)
\textbf{List}

\qquad

\textbf{পরিবর্তন করা যাবে না}
\(\longrightarrow\)
\textbf{Tuple}

\medskip

\textbf{Duplicate চাই না}
\(\longrightarrow\)
\textbf{Set}

\qquad

\textbf{Key দিয়ে Data চাই}
\(\longrightarrow\)
\textbf{Dictionary}

\end{center}

\end{rememberbox}

\small

%================================================
% 5.44 Function
%================================================

